\documentclass[12pt]{article}

\usepackage[margin=1in]{geometry}
\usepackage{enumitem}
\usepackage{amsmath}
\usepackage{amssymb}
\usepackage{array}
\usepackage{titlesec}

\title{Lecture 8 Notes: Aristotle's \textit{Topics}, Book I}
\author{}
\date{}

\begin{document}

\maketitle

\section*{Central Theme}

Aristotle's \textit{Topics}, Book I introduces dialectical reasoning: reasoning
not from scientific first principles, but from reputable or generally accepted
opinions.

\[
\boxed{
\textit{Topics}, \text{ Book I}
=
\text{Aristotle's introduction to dialectical reasoning}
}
\]

The \textit{Topics} teaches us how to argue responsibly from what people
already accept, especially when we are examining problems rather than
demonstrating scientific truths.

\section*{1. From Demonstration to Dialectic}

The \textit{Posterior Analytics} studied scientific demonstration.

\[
\textit{Posterior Analytics} = \text{scientific demonstration}
\]

The \textit{Topics} studies dialectical reasoning.

\[
\textit{Topics} = \text{dialectical reasoning}
\]

In demonstration, reasoning begins from true, primary, explanatory principles.
In dialectic, reasoning begins from generally accepted opinions.

\[
\boxed{
\text{Demonstration begins from first principles.}
}
\]

\[
\boxed{
\text{Dialectic begins from accepted opinions.}
}
\]

\subsection*{Teaching Point}

Aristotle is not abandoning logic. He is applying reason to a different kind of
setting: argument, inquiry, debate, training, and philosophical examination.

\[
\boxed{
\text{Dialectic works where inquiry has not yet reached demonstration.}
}
\]

\section*{2. The Program of the Treatise}

Aristotle says that the aim of the treatise is to find a line of inquiry by
which we can reason from generally accepted opinions about every problem set
before us.

It also aims to help us, when defending a position, avoid saying anything that
obstructs our own argument.

\subsection*{Whiteboard}

\[
\text{Find arguments from accepted opinions.}
\]

\[
\text{Avoid self-obstruction in debate.}
\]

\subsection*{Teaching Point}

The \textit{Topics} is a practical logical handbook.

\[
\boxed{
\text{It teaches how to find, organize, and use arguments.}
}
\]

\section*{3. What Is Reasoning?}

Aristotle defines reasoning generally.

\[
\boxed{
\text{Reasoning is an argument in which, certain things being laid down,}
}
\]

\[
\boxed{
\text{something other than these necessarily comes about through them.}
}
\]

This definition connects the \textit{Topics} back to the
\textit{Prior Analytics}. Reasoning is still structured necessity.

\subsection*{Teaching Point}

The difference between kinds of reasoning lies not in whether they reason, but
in the status of their premises.

\[
\boxed{
\text{The kind of reasoning depends on the kind of premises.}
}
\]

\section*{4. Four Kinds of Reasoning}

Aristotle distinguishes several kinds of reasoning.

\[
\begin{array}{ll}
1. & \text{Demonstrative reasoning} \\
2. & \text{Dialectical reasoning} \\
3. & \text{Contentious reasoning} \\
4. & \text{Mis-reasoning within a special science}
\end{array}
\]

\subsection*{Demonstrative Reasoning}

Demonstrative reasoning begins from premises that are true and primary, or from
premises known through what is true and primary.

\[
\boxed{
\text{Demonstration} \rightarrow \text{true and primary premises}
}
\]

\subsection*{Dialectical Reasoning}

Dialectical reasoning begins from generally accepted opinions.

\[
\boxed{
\text{Dialectic} \rightarrow \text{generally accepted opinions}
}
\]

\subsection*{Contentious Reasoning}

Contentious reasoning begins from opinions that merely seem to be generally
accepted, or merely seems to reason.

\[
\boxed{
\text{Contentious reasoning} \rightarrow \text{seeming acceptance or seeming inference}
}
\]

\subsection*{Mis-Reasoning in a Special Science}

Mis-reasoning may occur within a special science when someone reasons from
false assumptions that nevertheless belong to that science.

For example, in geometry one may draw a false figure or make an impossible
construction.

\[
\boxed{
\text{Scientific mis-reasoning} \rightarrow \text{false assumptions within a science}
}
\]

\section*{5. Generally Accepted Opinions}

Generally accepted opinions are those accepted by:

\[
\begin{array}{ll}
1. & \text{everyone} \\
2. & \text{the majority} \\
3. & \text{the philosophers} \\
4. & \text{all, most, or the most notable philosophers}
\end{array}
\]

These are often called \textit{endoxa}: reputable or generally accepted
opinions.

\[
\boxed{
\textit{Endoxa} = \text{reputable opinions}
}
\]

\subsection*{Teaching Point}

Aristotle does not mean that popular opinion is automatically true. He means
that dialectic begins from premises that an audience or interlocutor can
reasonably be expected to grant.

\[
\boxed{
\text{Dialectic tests opinions by reasoning from opinions.}
}
\]

\section*{6. Dialectic and Demonstration Compared}

\[
\begin{array}{c|c|c}
\textbf{Type} & \textbf{Premises} & \textbf{Aim} \\
\hline
\text{Demonstration} & \text{true and primary} & \text{scientific knowledge} \\
\text{Dialectic} & \text{generally accepted} & \text{examination and argument} \\
\text{Contentious reasoning} & \text{apparently accepted} & \text{victory or appearance} \\
\end{array}
\]

\subsection*{Teaching Point}

Dialectic is weaker than demonstration in certainty, but broader in use.

\[
\boxed{
\text{Dialectic is not science, but it is indispensable for inquiry.}
}
\]

\section*{7. The Three Uses of the \textit{Topics}}

Aristotle says the treatise is useful for three purposes.

\[
\begin{array}{ll}
1. & \text{Intellectual training} \\
2. & \text{Casual encounters} \\
3. & \text{The philosophical sciences}
\end{array}
\]

\subsection*{Intellectual Training}

Dialectic trains the mind by giving us a method for arguing about proposed
subjects.

\[
\boxed{
\text{Training} \rightarrow \text{practice in argument}
}
\]

\subsection*{Casual Encounters}

Dialectic helps in ordinary discussions. It allows us to meet people on the
ground of their own convictions and shift arguments that are unsound.

\[
\boxed{
\text{Conversation} \rightarrow \text{argue from another's beliefs}
}
\]

\subsection*{Philosophical Sciences}

Dialectic helps philosophy by enabling us to raise difficulties on both sides
of a question. This makes truth and error easier to detect.

It also helps with first principles, because the principles of each science
cannot be discussed from within that same science by means of its own
principles.

\[
\boxed{
\text{Philosophy} \rightarrow \text{test difficulties and examine principles}
}
\]

\subsection*{Teaching Point}

Dialectic is useful because inquiry often begins before demonstration is
possible.

\[
\boxed{
\text{Dialectic prepares the way toward principles.}
}
\]

\section*{8. Dialectic as Criticism}

Aristotle describes dialectic as a critical process.

It is especially useful in relation to the ultimate bases of the sciences,
because the first principles of a science cannot be proved from within that
science.

\subsection*{Example}

Geometry cannot prove its own first principles geometrically, because those
principles are prior to geometrical proof.

\subsection*{Teaching Point}

Dialectic provides a path toward the principles of inquiry.

\[
\boxed{
\text{Dialectic is the art of testing what demonstration must assume.}
}
\]

\section*{9. The Ideal of Dialectical Skill}

Aristotle compares dialectic to rhetoric and medicine.

A rhetorician cannot persuade everyone.

A doctor cannot heal everyone.

A dialectician cannot solve every argument in every situation.

But each has an art if he does not omit any available means.

\subsection*{Teaching Point}

Dialectical skill means using available argumentative materials well.

\[
\boxed{
\text{Dialectical skill} = \text{doing what one can with the available materials}
}
\]

\section*{10. Propositions and Problems}

Aristotle distinguishes propositions and problems.

Arguments begin from propositions.

Reasonings are about problems.

\[
\boxed{
\text{Proposition} = \text{claim offered for acceptance}
}
\]

\[
\boxed{
\text{Problem} = \text{question for inquiry}
}
\]

\subsection*{Example of a Proposition}

\[
\text{``Animal is the genus of man, is it not?''}
\]

\subsection*{Example of a Problem}

\[
\text{``Is animal the genus of man or not?''}
\]

\subsection*{Teaching Point}

The difference between a proposition and a problem is often the turn of phrase.

\[
\boxed{
\text{Dialectic moves between question-form and proposition-form.}
}
\]

\section*{11. The Four Predicables}

Every dialectical proposition and problem concerns one of four predicables.

\[
\begin{array}{ll}
1. & \text{Definition} \\
2. & \text{Property} \\
3. & \text{Genus} \\
4. & \text{Accident}
\end{array}
\]

These are called predicables because they classify how a predicate belongs to a
subject.

\[
\boxed{
\text{Every dialectical problem concerns definition, property, genus, or accident.}
}
\]

\subsection*{Teaching Point}

The four predicables give the basic map of dialectical problems.

\[
\boxed{
\text{The rest of the \textit{Topics} develops rules for arguing about these.}
}
\]

\section*{12. Definition}

A definition is a phrase signifying a thing's essence.

\[
\boxed{
\text{Definition} = \text{phrase signifying essence}
}
\]

A definition tells what a thing is.

\subsection*{Example}

\[
\text{man} = \text{animal that walks on two feet}
\]

\subsection*{Important Point}

A definition is not usually a single word. It is a phrase or formula.

\[
\boxed{
\text{A definition is always an account, not merely a name.}
}
\]

\subsection*{Teaching Point}

In the \textit{Topics}, Aristotle is concerned with how definitions function in
argument.

\[
\boxed{
\text{A definition can be tested dialectically.}
}
\]

\section*{13. Property}

A property is a predicate that belongs to a thing alone and is predicated
convertibly of it, but does not indicate its essence.

\[
\boxed{
\text{Property} =
\text{convertible predicate that does not state essence}
}
\]

\subsection*{Example}

It is a property of man to be capable of learning grammar.

\[
\text{man} \rightarrow \text{capable of learning grammar}
\]

\[
\text{capable of learning grammar} \rightarrow \text{man}
\]

So:

\[
\text{man} \leftrightarrow \text{capable of learning grammar}
\]

\subsection*{Teaching Point}

A property is exclusive, but not essential.

\[
\boxed{
\text{Property belongs only to the subject, but does not say what the subject is.}
}
\]

\section*{14. Genus}

A genus is predicated in the category of essence of several things differing in
kind.

\[
\boxed{
\text{Genus} = \text{what something is, said broadly of many species}
}
\]

\subsection*{Example}

If asked what a human being is, it is appropriate to answer:

\[
\text{animal}
\]

Animal is the genus of man.

\[
\text{man} \rightarrow \text{animal}
\]

\[
\text{ox} \rightarrow \text{animal}
\]

\subsection*{Teaching Point}

Genus belongs to essence, but it does not give the whole essence.

\[
\boxed{
\text{A genus tells what kind of thing something is, but not exactly what it is.}
}
\]

\section*{15. Accident}

An accident is what belongs to a thing but is neither definition, property, nor
genus.

Aristotle gives a second and clearer account: an accident is something that may
belong or not belong to one and the same thing.

\[
\boxed{
\text{Accident} = \text{what may belong or not belong to the same thing}
}
\]

\subsection*{Examples}

\[
\text{sitting}
\]

\[
\text{white}
\]

The same person may be sitting now and not sitting later.

The same thing may be white now and not white later.

\subsection*{Teaching Point}

Accident is predication without essential necessity.

\[
\boxed{
\text{An accident may be true of a subject without belonging to what it is.}
}
\]

\section*{16. Accident, Temporary Property, and Relative Property}

Aristotle notes that an accident can become a temporary or relative property.

\subsection*{Temporary Property}

If only one man is sitting, then sitting may function as a temporary property of
that man.

\[
\text{sitting} = \text{temporary property}
\]

\subsection*{Relative Property}

Two-footed may be treated as a property of man relative to horse and dog, even
if it is not a property absolutely.

\[
\text{two-footed} = \text{relative property}
\]

\subsection*{Teaching Point}

Dialectic must notice whether a predicate belongs absolutely, temporarily, or
relatively.

\[
\boxed{
\text{Many arguments fail by confusing absolute, temporary, and relative predication.}
}
\]

\section*{17. How the Four Predicables Relate}

\[
\begin{array}{c|c|c}
\textbf{Predicable} & \textbf{Convertible?} & \textbf{States Essence?} \\
\hline
\text{Definition} & \text{Yes} & \text{Yes} \\
\text{Property} & \text{Yes} & \text{No} \\
\text{Genus} & \text{No} & \text{Partially} \\
\text{Accident} & \text{No} & \text{No}
\end{array}
\]

\subsection*{Teaching Point}

The predicables classify the logical relation between a subject and predicate.

\[
\boxed{
\text{The question is not only what is said, but how it belongs.}
}
\]

\section*{18. Sameness}

Aristotle distinguishes several senses of sameness.

\[
\begin{array}{ll}
1. & \text{Numerical sameness} \\
2. & \text{Specific sameness} \\
3. & \text{Generic sameness}
\end{array}
\]

\subsection*{Numerical Sameness}

Numerical sameness occurs when there is one thing with more than one name.

\[
\text{doublet} = \text{cloak}
\]

\subsection*{Specific Sameness}

Specific sameness occurs when several things belong to the same species.

\[
\text{one man and another man}
\]

\subsection*{Generic Sameness}

Generic sameness occurs when several things belong to the same genus.

\[
\text{horse and man} \rightarrow \text{animal}
\]

\subsection*{Teaching Point}

Arguments about definition often depend on arguments about sameness and
difference.

\[
\boxed{
\text{To test a definition, ask whether the things identified are really the same.}
}
\]

\section*{19. Why the Four Predicables Are Exhaustive}

Aristotle argues that the four predicables cover all cases of predication.

Every predicate is either convertible with the subject or not.

\subsection*{If Convertible}

If the predicate is convertible with the subject, then it is either:

\[
\text{definition}
\]

or:

\[
\text{property}
\]

It is a definition if it signifies essence.

It is a property if it does not signify essence.

\subsection*{If Not Convertible}

If the predicate is not convertible with the subject, then it is either part of
the definition or not.

If it is part of the definition, it is genus or differentia.

If it is not part of the definition, it is accident.

\[
\begin{array}{c|c}
\textbf{Convertible} & \text{Definition or Property} \\
\hline
\textbf{Not Convertible} & \text{Genus/Differentia or Accident}
\end{array}
\]

\subsection*{Teaching Point}

The predicables classify all the basic ways a predicate may belong to a subject.

\[
\boxed{
\text{Every predicate has a dialectical status.}
}
\]

\section*{20. Predicables and Categories}

Aristotle connects the four predicables with the ten categories.

\[
\begin{array}{ll}
1. & \text{Essence} \\
2. & \text{Quantity} \\
3. & \text{Quality} \\
4. & \text{Relation} \\
5. & \text{Place} \\
6. & \text{Time} \\
7. & \text{Position} \\
8. & \text{State} \\
9. & \text{Activity} \\
10. & \text{Passivity}
\end{array}
\]

Accident, genus, property, and definition will always fall into one of these
categories.

\subsection*{Key Distinction}

\[
\text{Predicables} = \text{how a predicate belongs}
\]

\[
\text{Categories} = \text{what kind of predicate it is}
\]

\subsection*{Teaching Point}

Categories classify kinds of being. Predicables classify modes of predication.

\[
\boxed{
\text{Categories answer ``what kind?'' Predicables answer ``how does it belong?''}
}
\]

\section*{21. Dialectical Propositions}

A dialectical proposition is not just any proposition.

It is a question asking something held by all, most, or the philosophers,
provided that it is not contrary to general opinion.

\[
\boxed{
\text{Dialectical proposition} = \text{reputable claim suitable for argument}
}
\]

\subsection*{Examples}

\[
\text{Is the knowledge of contraries the same?}
\]

\[
\text{Should one do good to one's friends?}
\]

Dialectical propositions also include views like generally accepted opinions,
contradictories of contraries of generally accepted opinions, and claims in
accordance with recognized arts.

\subsection*{Teaching Point}

Dialectic begins from claims an interlocutor can reasonably be expected to
grant.

\[
\boxed{
\text{A dialectical premise must be reputable enough to be usable.}
}
\]

\section*{22. Dialectical Problems}

A dialectical problem is a subject of inquiry that contributes either to:

\[
\text{choice and avoidance}
\]

or to:

\[
\text{truth and knowledge}
\]

It may contribute directly, or as a help to some other problem.

\subsection*{Examples}

For choice and avoidance:

\[
\text{Is pleasure to be chosen or not?}
\]

For truth and knowledge:

\[
\text{Is the universe eternal or not?}
\]

\subsection*{Teaching Point}

A dialectical problem must be genuinely disputable and worth inquiry.

\[
\boxed{
\text{Dialectic studies questions where reasoned disagreement is possible.}
}
\]

\section*{23. Theses}

A thesis is a supposition of an eminent philosopher that conflicts with general
opinion.

\[
\boxed{
\text{Thesis} = \text{reputable philosophical claim contrary to common opinion}
}
\]

\subsection*{Examples}

\[
\text{Heraclitus: all things are in motion.}
\]

\[
\text{Melissus: Being is one.}
\]

A thesis is a kind of problem, but not every problem is a thesis.

\[
\text{Thesis} \subset \text{Problem}
\]

\subsection*{Teaching Point}

A thesis is worth examining because it is both serious and controversial.

\[
\boxed{
\text{Dialectic gives disciplined form to philosophical controversy.}
}
\]

\section*{24. What Should Not Be Examined Dialectically}

Aristotle says not every problem or thesis should be examined.

Some questions are morally corrupt rather than dialectically puzzling.

\[
\text{Should one honor the gods and love one's parents?}
\]

Such a person needs correction, not argument.

Other questions require perception rather than dialectic.

\[
\text{Is snow white?}
\]

Such a person needs to look, not debate.

\subsection*{Teaching Point}

Dialectic is for genuine puzzles, not for moral blindness or perceptual
ignorance.

\[
\boxed{
\text{Not every question deserves dialectical treatment.}
}
\]

\section*{25. Dialectical Reasoning and Induction}

Aristotle distinguishes two species of dialectical argument.

\[
\begin{array}{ll}
1. & \text{Reasoning} \\
2. & \text{Induction}
\end{array}
\]

\subsection*{Reasoning}

Reasoning proceeds from premises to a necessary conclusion.

\[
\text{premises} \rightarrow \text{necessary conclusion}
\]

\subsection*{Induction}

Induction passes from individuals to universals.

\[
\text{particulars} \rightarrow \text{universal}
\]

\subsection*{Example}

\[
\text{The skilled pilot is best at his task.}
\]

\[
\text{The skilled charioteer is best at his task.}
\]

\[
\therefore \text{In general, the skilled person is best at his task.}
\]

\subsection*{Teaching Point}

Induction is clearer and more persuasive to most people, while reasoning is
more forceful against contradictious people.

\[
\boxed{
\text{Dialectic uses both examples that persuade and arguments that compel.}
}
\]

\section*{26. The Four Sources of Arguments}

Aristotle explains how we become well supplied with arguments.

There are four sources.

\[
\begin{array}{ll}
1. & \text{Securing propositions} \\
2. & \text{Distinguishing the senses of terms} \\
3. & \text{Discovering differences} \\
4. & \text{Investigating likenesses}
\end{array}
\]

\subsection*{Teaching Point}

To argue well, one must prepare materials before the argument begins.

\[
\boxed{
\text{Dialectic requires a supply of propositions, meanings, differences, and likenesses.}
}
\]

\section*{27. Securing Propositions}

Aristotle recommends collecting propositions from several sources.

\[
\begin{array}{ll}
1. & \text{opinions held by all} \\
2. & \text{opinions held by most} \\
3. & \text{opinions held by philosophers} \\
4. & \text{opinions contrary to what seems generally accepted} \\
5. & \text{opinions in accordance with recognized arts} \\
6. & \text{opinions like those generally accepted}
\end{array}
\]

\subsection*{Practical Method}

Draw up lists under headings.

\[
\text{On Good}
\]

\[
\text{On Life}
\]

\[
\text{On Pleasure}
\]

\[
\text{On Justice}
\]

Also note the views of individual thinkers.

\[
\text{Empedocles says the elements of bodies are four.}
\]

\subsection*{Teaching Point}

Dialectic requires preparation: collect and organize argumentative material.

\[
\boxed{
\text{The dialectician is well supplied because he has already gathered endoxa.}
}
\]

\section*{28. Ethical, Natural, and Logical Propositions}

Aristotle roughly divides propositions and problems into three kinds.

\[
\begin{array}{ll}
1. & \text{Ethical} \\
2. & \text{Natural-philosophical} \\
3. & \text{Logical}
\end{array}
\]

\subsection*{Ethical Example}

\[
\text{Should one obey one's parents or the laws, if they disagree?}
\]

\subsection*{Natural-Philosophical Example}

\[
\text{Is the universe eternal or not?}
\]

\subsection*{Logical Example}

\[
\text{Is the knowledge of opposites the same or not?}
\]

\subsection*{Teaching Point}

Dialectic can discuss many subjects because it works from reputable opinion
rather than from the proper principles of a single science.

\[
\boxed{
\text{Dialectic ranges widely because its starting points are common opinions.}
}
\]

\section*{29. Philosophy and Dialectic Treat Propositions Differently}

For philosophy, propositions should be treated according to truth.

For dialectic, propositions are treated with an eye to general opinion.

\[
\text{Philosophy} \rightarrow \text{truth}
\]

\[
\text{Dialectic} \rightarrow \text{general opinion}
\]

\subsection*{Teaching Point}

Dialectic is not indifferent to truth, but its immediate method is to reason
from what is accepted.

\[
\boxed{
\text{Dialectic is truth-seeking through reputable opinion.}
}
\]

\section*{30. Distinguishing Ambiguous Meanings}

Aristotle says we must distinguish how many senses a term bears.

It is not enough to list the senses. We should also try to define them.

\subsection*{Example: Good}

The term ``good'' may be said in different senses.

\[
\text{virtue is good}
\]

\[
\text{what produces health is good}
\]

\[
\text{what produces vigor is good}
\]

In one case, good is said because of an intrinsic quality. In another, it is
said because something produces a result.

\subsection*{Teaching Point}

One word may hide many questions.

\[
\boxed{
\text{Clarifying meanings prevents arguments from being aimed merely at words.}
}
\]

\section*{31. Tests for Ambiguity}

Aristotle gives many ways to test whether a term has more than one meaning.

\[
\begin{array}{ll}
1. & \text{Examine its contraries} \\
2. & \text{Examine intermediates} \\
3. & \text{Examine contradictory opposites} \\
4. & \text{Examine privation and possession} \\
5. & \text{Examine inflected forms} \\
6. & \text{Examine categories} \\
7. & \text{Examine genera} \\
8. & \text{Examine definitions in combination} \\
9. & \text{Ask whether it admits of more and less} \\
10. & \text{Examine differentiae}
\end{array}
\]

\subsection*{Teaching Point}

Ambiguity is discovered by testing the relations around a term.

\[
\boxed{
\text{A word's meaning is revealed by its logical neighborhood.}
}
\]

\section*{32. Ambiguity Through Contraries}

A term may be ambiguous if its contrary differs in different contexts.

\subsection*{Example}

``Sharp'' has different contraries.

\[
\text{sharp note} \leftrightarrow \text{flat note}
\]

\[
\text{sharp edge} \leftrightarrow \text{dull edge}
\]

Therefore ``sharp'' is not being used in the same sense in both cases.

\subsection*{Teaching Point}

Different contraries often reveal different meanings.

\[
\boxed{
\text{If the opposites differ, the original term may be ambiguous.}
}
\]

\section*{33. Ambiguity Through Categories}

A term may be ambiguous if it signifies different categories in different uses.

\subsection*{Example}

``Good'' may signify:

\[
\text{quality, as in a virtuous soul}
\]

\[
\text{relation to health, as in good medicine}
\]

\[
\text{time, as in the right occasion}
\]

\[
\text{quantity, as in the proper amount}
\]

\subsection*{Teaching Point}

If the same word ranges across categories, it may not have one meaning.

\[
\boxed{
\text{Category shifts often signal ambiguity.}
}
\]

\section*{34. Ambiguity Through Genera}

A term may be ambiguous if the things named by it belong to different genera
that are not subaltern.

\subsection*{Example}

``Donkey'' may refer to:

\[
\text{an animal}
\]

or:

\[
\text{an engine}
\]

Since these belong to different genera, the term is ambiguous.

\subsection*{Teaching Point}

Different non-subaltern genera reveal different definitions.

\[
\boxed{
\text{If the genera differ, the meaning probably differs.}
}
\]

\section*{35. Why Ambiguity Matters}

Distinguishing meanings is useful for several reasons.

\[
\begin{array}{ll}
1. & \text{It promotes clarity.} \\
2. & \text{It prevents reasoning merely about words.} \\
3. & \text{It prevents being misled by false reasoning.} \\
4. & \text{It helps expose an opponent's equivocation.}
\end{array}
\]

\subsection*{Teaching Point}

Dialectic requires clarity about what is being discussed.

\[
\boxed{
\text{If the meaning is unclear, the argument may not have a single subject.}
}
\]

\section*{36. Differences}

Aristotle says we should examine the differences between things.

Differences should especially be examined within the same genus.

\subsection*{Examples Within the Same Genus}

\[
\text{How does justice differ from courage?}
\]

\[
\text{How does wisdom differ from temperance?}
\]

These all belong to the genus of virtue.

\subsection*{Example Across Genera}

We may also compare across genera if they are not too far apart.

\[
\text{How does sensation differ from knowledge?}
\]

\subsection*{Teaching Point}

To know what a thing is, know how it differs from its neighbors.

\[
\boxed{
\text{Difference helps with sameness, definition, and classification.}
}
\]

\section*{37. Likenesses}

Aristotle also says we should investigate likenesses.

Likeness is especially useful among things belonging to different genera.

\subsection*{Analogical Form}

\[
A:B = C:D
\]

\subsection*{Example}

\[
\text{knowledge} : \text{object of knowledge}
=
\text{sensation} : \text{object of sensation}
\]

Another example:

\[
\text{sight is in the eye as reason is in the soul}
\]

\subsection*{Teaching Point}

Likeness lets us transfer insight from one case to another.

\[
\boxed{
\text{Analogy supplies arguments across different domains.}
}
\]

\section*{38. Uses of Likeness}

The examination of likeness is useful for:

\[
\begin{array}{ll}
1. & \text{inductive arguments} \\
2. & \text{hypothetical reasonings} \\
3. & \text{rendering definitions}
\end{array}
\]

\subsection*{Induction}

Induction depends on recognizing likeness among particular cases.

\[
\text{similar cases} \rightarrow \text{universal claim}
\]

\subsection*{Hypothetical Reasoning}

If it is generally accepted that what is true of one similar case is true of the
others, then we may use likeness to reason by hypothesis.

\subsection*{Definition}

Likeness helps us see the common genus or common feature.

\[
\text{point} : \text{line}
=
\text{unit} : \text{number}
\]

Each is a kind of starting point.

\subsection*{Teaching Point}

Likeness is not merely decorative. It is a tool for finding universals,
hypotheses, and definitions.

\[
\boxed{
\text{Similarity helps the dialectician move from case to principle.}
}
\]

\section*{39. The Architecture of Book I}

Book I introduces the whole project of the \textit{Topics}.

\[
\text{program of dialectic}
\rightarrow
\text{kinds of reasoning}
\rightarrow
\text{uses of dialectic}
\rightarrow
\text{propositions and problems}
\]

\[
\rightarrow
\text{four predicables}
\rightarrow
\text{sameness}
\rightarrow
\text{categories}
\rightarrow
\text{dialectical propositions and problems}
\]

\[
\rightarrow
\text{reasoning and induction}
\rightarrow
\text{four sources of arguments}
\]

\subsection*{Teaching Point}

Book I is not yet the full collection of topical rules. It is the introduction
that explains what dialectic is, what it is for, what its materials are, and how
arguments are supplied.

\[
\boxed{
\text{Book I gives the map of dialectical inquiry.}
}
\]

\section*{40. Summary of Main Ideas}

\begin{enumerate}[itemsep=0.5em]
    \item The \textit{Topics} studies dialectical reasoning from generally
    accepted opinions.

    \item Reasoning is an argument in which something follows necessarily from
    what has been laid down.

    \item Demonstrative reasoning begins from true and primary premises.

    \item Dialectical reasoning begins from reputable opinions.

    \item Contentious reasoning begins from merely apparent reputable opinions
    or merely appears to reason.

    \item Dialectic is useful for intellectual training, casual encounters, and
    philosophical inquiry.

    \item Dialectic is especially important for examining first principles.

    \item Propositions provide the materials of argument; problems provide the
    subjects of inquiry.

    \item Every dialectical problem concerns definition, property, genus, or
    accident.

    \item Definition signifies essence.

    \item Property is convertible with the subject but does not signify essence.

    \item Genus is predicated in the category of essence of things differing in
    kind.

    \item Accident may belong or not belong to the same thing.

    \item The predicables classify how predicates belong; the categories
    classify what kind of predicate is involved.

    \item Dialectical propositions and problems must be reputable, disputable,
    and useful.

    \item Induction moves from particulars to universals; reasoning moves from
    premises to necessary conclusion.

    \item The four sources of arguments are propositions, meanings,
    differences, and likenesses.

    \item Ambiguity must be carefully detected and resolved.

    \item Differences help define and distinguish things.

    \item Likenesses support induction, hypothesis, and definition.
\end{enumerate}

\section*{Final Framing}

\[
\boxed{
\textit{Topics}, \text{ Book I, teaches the art of finding arguments from reputable opinions.}
}
\]

The \textit{Posterior Analytics} showed how science demonstrates from first
principles. The \textit{Topics} shows how philosophy trains itself toward
principles by testing, organizing, and arguing from what is generally accepted.

\[
\boxed{
\text{Dialectic is disciplined reasoning in the space between opinion and science.}
}
\]

\end{document}